\section{Lastenheft: Web-App Turnierplaner}

\subsection{Zielsetzung}
Das Projekt umfasst die Entwicklung einer webbasierten Applikation zur Organisation und Verwaltung von Sportturnieren. Ziel ist die Automatisierung von Planungs- und Durchführungsprozessen (Baum-Generierung) sowie die Echtzeit-Visualisierung des Turnierverlaufs.

\subsection{Funktionale Anforderungen (FA)}

\textbf{Team-Management}
\begin{itemize}
    \item Erfassung von Teamnamen und Spielerlisten (CRUD).
    \item Bearbeitung und Löschung von Teams vor Turnierstart.
\end{itemize}

\textbf{Turniermodus \& Baum-Generierung}
\begin{itemize}
    \item Unterstützung von K.-o.-System (Single Elimination)
    \item Automatisierte Generierung des Turnierbaums (zufällig Setzliste).
\end{itemize}

\textbf{Dashboard \& Monitoring}
\begin{itemize}
    \item Zentrales Dashboard zur Anzeige des aktuellen Turnierstands.
    \item Automatische Ermittlung des Gewinners und Update des Turnierbaums.
\end{itemize}

\subsection{Nicht-funktionale Anforderungen (NFA)}
\begin{itemize}
    \item \textbf{Zuverlässigkeit:} Persistente Speicherung der Turnierdaten zur Absicherung bei Verbindungsabbruch.
\end{itemize}

\subsection{Logik-Parameter}
Für die Berechnung der Turnierstruktur gelten folgende mathematische Grundlagen:

\begin{itemize}
    \item \textbf{Anzahl der Spiele ($N$):} 
    \[ N = n - 1 \] 
    (wobei $n$ die Anzahl der Teams im K.-o.-System ist).
    
    \item \textbf{Rundenanzahl ($R$):} 
    \[ R = \lceil \log_2(n) \rceil \]
\end{itemize}